\pdfoutput=1
\documentclass[11pt]{article}

\usepackage[T1]{fontenc}
\usepackage{lmodern}
\usepackage{amsmath,amssymb}
\usepackage{booktabs}
\usepackage{enumitem}
\usepackage[margin=1in]{geometry}
\usepackage{hyperref}
\usepackage{xcolor}
\usepackage{listings}

\hypersetup{
  colorlinks=true,
  linkcolor=blue!55!black,
  citecolor=blue!55!black,
  urlcolor=blue!55!black
}

\lstset{
  basicstyle=\ttfamily\small,
  columns=fullflexible,
  breaklines=true,
  frame=single,
  xleftmargin=0.5em,
  xrightmargin=0.5em
}

\newcommand{\qedesk}{\textsc{QEDesk}}
\newcommand{\lean}{Lean 4}
\newcommand{\mathlib}{Mathlib}
\newcommand{\mcp}{MCP}

\title{\qedesk: A Local Proof-Audit Workstation for Undergraduate Mathematics}

\author{
  Dongwoo Lee\\
  \small QEDesk project\\
  \small \url{https://github.com/dev-heps/QEDesk}
}

\date{Draft version: July 7, 2026}

\begin{document}
\maketitle

\begin{abstract}
\qedesk{} is a local workstation for undergraduate mathematics that combines
\lean{}, LaTeX, Python, OpenCode-style agent workflows, OpenRouter-compatible
model calls, and model-context-protocol tooling.  Its central design rule is
that language models may audit, decompose, explain, and propose candidates, but
\lean{} remains the final verifier.  The system is intended for students who
write informal proofs in LaTeX and want structured feedback about missing
justifications, quantifier mistakes, hidden hypotheses, set-membership errors,
and statement drift before attempting a full formalization.  This technical
report describes the prototype architecture, the proof-audit data contracts,
the blueprint-and-DAG workflow, and a budget-aware model-routing layer.  We
position \qedesk{} as an educational proof-audit workstation rather than an
autonomous theorem-proving benchmark system.
\end{abstract}

\section{Introduction}

Students often learn proof writing through a fragile loop: write an informal
argument, receive delayed feedback, and then discover that a step described as
``obvious'' hides a missing hypothesis or a quantifier error.  Large language
models can give immediate feedback, but they can also smooth over precisely the
logical gaps that the student needs to see.  At the same time, theorem provers
such as \lean{} can certify formal proofs, but the jump from informal prose to
fully formal proof scripts is steep for beginners.

\qedesk{} is designed as a local bridge between these two modes.  The student
keeps a human-readable LaTeX notebook, writes or sketches Lean targets in a
nearby file, and uses an assistant only as a proof auditor and candidate
generator.  The assistant is not allowed to declare success merely because a
response sounds plausible.  A mathematical claim is considered formally checked
only after \lean{} accepts it.

The project is deliberately modest.  It is not a new theorem prover, not a
replacement for \mathlib{}, and not a benchmark claim that one particular model
is the best at mathematics.  Its contribution is an engineering pattern for
local, student-facing proof feedback: decompose first, audit locally, formalize
one node at a time, and use the prover as the authority.

This report makes three concrete contributions:

\begin{enumerate}[leftmargin=2em]
  \item a proof-audit protocol that classifies local gaps in informal
  undergraduate proofs before asking for Lean tactics;
  \item an artifact structure that connects LaTeX notes, Lean declarations,
  Lean Blueprint pages, and a QEDesk-specific dependency sidecar; and
  \item a budget-aware model-routing policy that records cost and treats model
  output as advisory until Lean verifies the relevant formal target.
\end{enumerate}

\section{Background and Related Work}

\paragraph{Lean and Mathlib.}
\lean{} is both a theorem prover and a programming language with an extensible
elaborator, tactic framework, and editor-facing language server
\cite{lean4}.  \mathlib{} provides the large mathematical library on which many
Lean developments depend, together with a community process for maintaining
formalized mathematics at scale \cite{mathlib}.  \qedesk{} builds on this
ecosystem rather than replacing it: Lean provides the formal kernel and
\mathlib{} supplies the mathematical vocabulary.

\paragraph{Open theorem-proving environments.}
Recent systems such as LeanDojo expose Lean programmatically for data
extraction, premise selection, environment interaction, and reproducible
machine-learning experiments \cite{leandojo}.  REAL-Prover is another
retrieval-augmented Lean 4 prover that focuses on stepwise theorem proving and
college-level mathematical reasoning benchmarks \cite{realprover}.  These
systems motivate the use of retrieval, local proof states, and iterative Lean
feedback.  \qedesk{} differs in emphasis: it targets a study workstation whose
default task is to audit a student's informal reasoning, not to maximize
Pass@\(k\) on a prover benchmark.

\paragraph{Education-oriented Lean interfaces.}
Verbose Lean uses controlled natural language and custom Lean tooling to help
undergraduate students transfer proof-assistant work back to traditional paper
proofs \cite{verboselean}.  \qedesk{} shares this pedagogical concern but keeps
the student's primary artifact as LaTeX prose.  Rather than replacing proof
writing with a controlled language, it audits the prose, records a blueprint,
and asks Lean to check only stable local targets.

\paragraph{Blueprints, automation, and budgeted model use.}
Lean Blueprint provides a way to organize formalization projects around
definitions, lemmas, theorem statements, and dependency graphs
\cite{leanblueprint}.  Automated tactics such as Aesop show how bounded proof
search can be useful when it remains inside Lean's checking discipline
\cite{aesop}.  On the model side, cascade-style approaches such as FrugalGPT
motivate routing routine calls to cheaper models and reserving expensive calls
for harder cases \cite{frugalgpt}.  \qedesk{} combines these ideas in a
student-facing workflow: graph first, local proof state second, model hint
third, and Lean verification last.

\section{Design Requirements}

\qedesk{} follows four design requirements.

\begin{enumerate}[leftmargin=2em]
  \item \textbf{Local first.}  The workstation keeps the student's LaTeX notes,
  Lean files, generated graphs, and cost ledger in the project directory or in
  explicitly disposable local build outputs.

  \item \textbf{Verifier first.}  The assistant can propose a repair, a lemma,
  or a tactic family, but \lean{} is the only component allowed to verify a
  formal proof.

  \item \textbf{Pedagogical hints before solutions.}  The default response is a
  local hint or gap classification, not a polished complete proof.

  \item \textbf{Replaceable models.}  Model names and prices are configuration
  details.  The workflow should continue to make sense when the cheap model,
  strong model, or provider changes.
\end{enumerate}

\section{Architecture}

The workstation separates retrieval, proof-state communication, formal
verification, and language-model reasoning.  Table~\ref{tab:components}
summarizes the boundary between components.

\begin{table}[h]
\centering
\small
\begin{tabular}{lll}
\toprule
Component & Responsibility & Non-responsibility \\
\midrule
RAG/search & Retrieve definitions, lemmas, notes, examples & Decide correctness \\
\mcp{} & Expose current Lean goals and diagnostics & Store global knowledge \\
\lean{} & Check formal statements and proof candidates & Teach prose style \\
Language model & Audit, decompose, explain, propose & Certify proofs \\
LaTeX & Preserve the student's notebook & Serve as formal truth \\
\bottomrule
\end{tabular}
\caption{QEDesk component boundaries.}
\label{tab:components}
\end{table}

The intended loop is local rather than monolithic:

\begin{lstlisting}
LaTeX proof note
  -> local proof audit
  -> blueprint DAG
  -> one Lean statement or lemma
  -> Lean check through Lake or MCP
  -> repair hint
  -> updated student note
\end{lstlisting}

This loop avoids sending the entire project, entire chat transcript, or full
\mathlib{} context to a model on every request.  It also avoids asking a model
to produce a complete proof script before the statement and local dependencies
are stable.

\section{Proof-Audit Protocol}

The proof-audit protocol begins by naming the assumptions, variables,
quantifiers, and exact goal.  The proof is then split into local inference
steps.  Each nontrivial step is marked as one of:

\begin{itemize}[leftmargin=2em]
  \item \texttt{verified};
  \item \texttt{missing-justification};
  \item \texttt{false};
  \item \texttt{unclear}.
\end{itemize}

For implementation, \qedesk{} uses a fixed gap taxonomy:

\begin{center}
\small
\begin{tabular}{lll}
\texttt{quantifier-jump} & \texttt{missing-hypothesis} & \texttt{missing-justification} \\
\texttt{false-inference} & \texttt{undefined-symbol} & \texttt{ambiguous-expression} \\
\texttt{set-membership-gap} & \texttt{calculation-gap} & \texttt{visual-dependence} \\
\texttt{statement-drift} & \texttt{unspecified} & \\
\end{tabular}
\end{center}

Hint levels range from level 0, which only highlights the suspicious span, to
level 4, which may include a Lean candidate patch.  Level 4 output remains a
candidate until Lean checks it.

The protocol is intentionally resistant to compressed answers.  Phrases such
as ``clearly'' or ``trivially'' are treated as insufficient unless they are
backed by a named definition, theorem, calculation, or Lean-checkable argument.

\section{Blueprint and DAG Workflow}

For nontrivial exercises, \qedesk{} first builds a small blueprint.  A blueprint
contains definitions, local lemmas, the final theorem, and dependency edges.
Each node may have an informal statement, an optional Lean declaration name,
a status, and gap metadata.  In the current prototype, the editable source of
this graph is the student's LaTeX file, using comments that do not affect PDF
rendering:

\begin{lstlisting}[language=TeX]
% QEDesk[node=mem_power_set][kind=lemma]
% QEDesk[uses=def_power_set]
\end{lstlisting}

The command-line synchronizer reads these comments and the Lean declarations,
then generates three disposable artifacts:

\begin{itemize}[leftmargin=2em]
  \item Lean Blueprint source under \texttt{blueprint/src/content.tex};
  \item a machine-readable sidecar at \texttt{build/qedesk-dag.json};
  \item a Mermaid dependency graph at \texttt{build/qedesk-dag.mmd}.
\end{itemize}

The web view is generated by Lean Blueprint.  \qedesk{} adds only a small view
switching layer around the generated pages, while treating the upstream
Blueprint renderer as the graph and document viewer.

\section{Token-Aware Model Routing}

\qedesk{} treats model calls as a budgeted resource.  The default cascade is a
cheap worker, then a stronger planner, then a deeper critic.  In the current
configuration these roles are mapped to OpenRouter-compatible DeepSeek model
identifiers, but the architecture is not tied to those names.

\begin{table}[h]
\centering
\small
\begin{tabular}{llll}
\toprule
Stage & Primary role & Escalation role & Context sent \\
\midrule
Audit & cheap worker & strong planner & local node and snippet \\
Blueprint & cheap worker & strong planner & node plus parents \\
Formalize & cheap worker & strong planner & statement skeleton \\
Repair & cheap worker & critic & Lean error excerpt \\
Verify & Lean server & none & formal file \\
\bottomrule
\end{tabular}
\caption{Budget-aware routing policy.}
\label{tab:routing}
\end{table}

The router prototype supports dry runs, real OpenRouter calls, and a SQLite
cost ledger.  It logs the worksheet id, node id, stage, model name, prompt and
completion token counts, estimated cost, elapsed time, outcome, and Lean
status.  Its output is advisory JSON; it is not proof verification.

\section{Implementation}

The repository is organized around a small set of files:

\begin{lstlisting}
src/main.tex           student-facing LaTeX note
src/Proof.lean         Lean statements and proofs
AGENTS.md              assistant behavior rules
docs/ARCHITECTURE.md   system design
docs/DATA_CONTRACTS.md gap labels and DAG contracts
tools/qedesk_sync.py   LaTeX/Lean/Blueprint synchronizer
tools/qedesk_router.py OpenRouter routing prototype
bin/qedesk             command-line wrapper
\end{lstlisting}

The common local workflow is:

\begin{lstlisting}
./bin/qedesk audit --dry-run --max-nodes 1 src/main.tex
./bin/qedesk sync
./bin/qedesk lean
./bin/qedesk pdf
./bin/qedesk blueprint
\end{lstlisting}

The Docker-based workstation can provide \lean{}, Lake, LaTeX, Python,
\texttt{lean-lsp-mcp}, Lean Blueprint, and related tools in one environment.
For source-only review, the repository can also be cloned and inspected without
running the container.  This distinction is important for distribution:
generated artifacts, cost ledgers, local API keys, Lean build products, and
Docker images are not part of the source archive.

\section{Current Status and Evaluation Plan}

\qedesk{} is currently a prototype artifact.  The implemented parts are the
source layout, command wrappers, documentation, proof-audit contracts, Lean
Blueprint synchronization, Mermaid sidecar generation, and a conservative
OpenRouter routing prototype.  The system does not yet include a controlled
user study, a large benchmark evaluation, or learned routing thresholds.

For a tool paper of this kind, the relevant near-term evaluation is not a
single theorem-proving score.  We instead track four artifact-level questions:

\begin{enumerate}[leftmargin=2em]
  \item \textbf{Faithfulness.}  Does the Lean statement still match the
  student's informal claim, or has statement drift occurred?
  \item \textbf{Locality.}  Does the assistant repair the failing node rather
  than rewriting the whole proof?
  \item \textbf{Verification.}  Does every claimed formal success correspond to
  a successful Lean check?
  \item \textbf{Cost.}  How many model tokens and retries are spent per useful
  proof-audit iteration?
\end{enumerate}

Future quantitative evaluation should therefore report cost per checked node,
number of repaired local gaps, statement-drift incidents, and Lean-verified
successes.  Until such an evaluation exists, \qedesk{} should be read as a
reproducible workflow proposal and open-source prototype rather than as a
state-of-the-art prover.

\section{Example Use}

Consider a typical set-theoretic statement about a function
\(f : A \to B\) and a subset \(B_1 \subseteq B\):

\[
  f(f^{-1}(B_1)) \subseteq B_1.
\]

A compressed informal proof may jump from
\(b \in f(f^{-1}(B_1))\) to \(b \in B_1\) without expanding the definition of
image and preimage.  \qedesk{} should not immediately write the finished
solution.  Instead, it should expose the local obligation:

\begin{quote}
If \(b \in f(f^{-1}(B_1))\), then there exists an \(a\) such that
\(a \in f^{-1}(B_1)\) and \(f(a)=b\).  To conclude \(b \in B_1\), expand the
definition of preimage to obtain \(f(a) \in B_1\), then rewrite using
\(f(a)=b\).
\end{quote}

At a higher hint level, the corresponding Lean node can be attempted and
checked.  If Lean rejects the candidate, the repair prompt should include only
the local goal state or concise error excerpt, not the full worksheet.

\section{Limitations}

The current prototype is not a complete autonomous proof agent.  Several parts
are intentionally conservative:

\begin{itemize}[leftmargin=2em]
  \item The synchronizer relies on explicit comments rather than unrestricted
  natural-language parsing.
  \item Router output is structured advice, not formal evidence.
  \item MCP integration depends on the editor and client environment.
  \item The system does not claim benchmark-level theorem-proving performance.
  \item Cost estimates depend on provider prices and must be refreshed.
\end{itemize}

These limitations are acceptable for the intended educational use case.  The
goal is to make logical gaps visible and to keep the student close to the
verification loop, not to hide the work behind an automatic final answer.

\section{Artifact Availability and Reproducibility}

The project is released as open-source software under the MIT license:

\begin{center}
\url{https://github.com/dev-heps/QEDesk}
\end{center}

The repository includes setup instructions, command wrappers, data contracts,
and documentation for the proof-audit workflow.  API keys, local cost ledgers,
intermediate LaTeX build products, Lean build products, and Docker images are
excluded from version control.  The compiled draft of this report is checked in
as a convenience artifact.

The intended source-level reproduction path is:

\begin{lstlisting}
git clone https://github.com/dev-heps/QEDesk
cd QEDesk
./bin/qedesk audit --dry-run --max-nodes 1 src/main.tex
\end{lstlisting}

The optional full workstation path additionally starts the container and builds
Lean and LaTeX artifacts:

\begin{lstlisting}
./bin/qedesk start
./bin/qedesk prepare
./bin/qedesk lean
./bin/qedesk pdf
./bin/qedesk blueprint
\end{lstlisting}

Real model calls require a user-provided OpenRouter key in a local
\texttt{.env} file.  This key is deliberately not part of the archived source.

\section{Conclusion}

\qedesk{} frames AI-assisted mathematics as a disciplined verification loop.
The student writes informal mathematics, the model audits and proposes local
next steps, the blueprint records dependencies, and \lean{} decides which
formal candidates actually pass.  This architecture shifts the engineering
focus away from finding a single perfect model and toward building a workflow
that remains useful as models, prices, and theorem-proving tools change.

\begin{thebibliography}{99}

\bibitem{lean4}
Leonardo de Moura and Sebastian Ullrich.
\newblock The Lean 4 theorem prover and programming language.
\newblock In \emph{Automated Deduction -- CADE 28}, 2021.

\bibitem{mathlib}
The mathlib community.
\newblock The Lean mathematical library.
\newblock In \emph{Proceedings of the 9th ACM SIGPLAN International
Conference on Certified Programs and Proofs}, 2020.

\bibitem{leanblueprint}
Patrick Massot.
\newblock Lean Blueprint.
\newblock \url{https://github.com/PatrickMassot/leanblueprint}.

\bibitem{leandojo}
Kaiyu Yang, Aidan M. Swope, Alex Gu, Rahul Chalamala, Peiyang Song,
Shixing Yu, Saad Godil, Ryan Prenger, and Anima Anandkumar.
\newblock LeanDojo: Theorem proving with retrieval-augmented language models.
\newblock \emph{arXiv:2306.15626}, 2023.

\bibitem{realprover}
Ziju Shen, Naohao Huang, Fanyi Yang, Yutong Wang, Guoxiong Gao, Tianyi Xu,
Jiedong Jiang, Wanyi He, Pu Yang, Mengzhou Sun, Haocheng Ju, Peihao Wu,
Bryan Dai, and Bin Dong.
\newblock REAL-Prover: Retrieval augmented Lean prover for mathematical
reasoning.
\newblock \emph{arXiv:2505.20613}, 2025.

\bibitem{verboselean}
Patrick Massot.
\newblock Teaching mathematics using Lean and controlled natural language.
\newblock In \emph{15th International Conference on Interactive Theorem
Proving (ITP 2024)}, 2024.

\bibitem{aesop}
Jannis Limperg and Asta Halkj{\ae}r From.
\newblock Aesop: White-box best-first proof search for Lean.
\newblock In \emph{Proceedings of the 12th ACM SIGPLAN International
Conference on Certified Programs and Proofs}, 2023.

\bibitem{frugalgpt}
Lingjiao Chen, Matei Zaharia, and James Zou.
\newblock FrugalGPT: How to use large language models while reducing cost and
improving performance.
\newblock \emph{arXiv:2305.05176}, 2023.

\bibitem{modelcontextprotocol}
Model Context Protocol.
\newblock \url{https://modelcontextprotocol.io/}.

\bibitem{qedesk}
QEDesk contributors.
\newblock QEDesk source repository.
\newblock \url{https://github.com/dev-heps/QEDesk}.

\end{thebibliography}

\end{document}
